\chapter{Transferaufgabe: Eigenes Datenprojekt}
\label{chap:transfer}

\begin{myremark}[Algorithmusart in dieser DL]
Funktionsart: \textbf{anwendungsübergreifend} (Klassifikation, Regression, Assoziation oder Filterung). \\
Umsetzungsart: \textbf{wahlweise regelbasiert oder selbstlernend}.
\end{myremark}

\begin{myremark}
Dieses Kapitel bildet den Abschluss des Skripts.
Sie setzt Kapitel~\cref{chap:einfuehrung} bis Kapitel~\cref{chap:randomforest} (und idealerweise Kapitel~\cref{chap:modellvergleich}) voraus.
Ziel ist die selbstständige Anwendung des Gelernten auf einen eigenen Datensatz.
\end{myremark}

\section{Projektauftrag}

\begin{myexercise}[Mini-Datenprojekt]
Entwerfen und realisieren Sie ein eigenes Datenprojekt nach folgendem Raster:

\begin{enumerate}
    \item \textbf{Problemdefinition:}
    \begin{itemize}
        \item Was soll vorhergesagt oder klassifiziert werden?
        \item Handelt es sich um Klassifikation oder Regression?
    \end{itemize}

    \item \textbf{Datenbeschaffung:}
    \begin{itemize}
        \item Mindestens \textbf{50 Zeilen} mit beschrifteten Daten
        \item Quelle dokumentieren (z.\,B. Kaggle, UCI ML Repository, eigene Erhebung)
        \item Lizenzbedingungen beachten
    \end{itemize}

    \item \textbf{Datenvorbereitung:}
    \begin{itemize}
        \item Fehlende Werte behandeln
        \item Features normalisieren falls nötig
        \item Train-Test-Split (80/20)
    \end{itemize}

    \item \textbf{Modellierung:}
    \begin{itemize}
        \item Mindestens \textbf{zwei} Modelle trainieren und vergleichen
        \item Begründung der Modellwahl
        \item Gütemasse berechnen (Accuracy/F1 oder MAE/RMSE)
    \end{itemize}

    \item \textbf{Visualisierung:}
    \begin{itemize}
        \item Mindestens ein aussagekräftiger Plot (z.\,B. Konfusionsmatrix, Regressionslinie,
              Feature Importance)
    \end{itemize}

    \item \textbf{Reflexion:}
    \begin{itemize}
        \item Was funktioniert gut, was nicht?
        \item Mögliche Bias-Risiken?
        \item Was würden Sie bei mehr Zeit verbessern?
    \end{itemize}
\end{enumerate}
\end{myexercise}

\section{Bewertungskriterien}

\begin{table}[H]
\centering
\small
\begin{tabular}{p{5.5cm}p{1.5cm}}
\toprule
\textbf{Kriterium} & \textbf{Punkte} \\
\midrule
Problemdefinition klar und sinnvoll & 1 \\
Datensatz geeignet und dokumentiert & 2 \\
Datenvorbereitung korrekt & 2 \\
Zwei Modelle korrekt trainiert und verglichen & 3 \\
Metriken korrekt berechnet und interpretiert & 2 \\
Visualisierung aussagekräftig & 2 \\
Reflexion kritisch und substanziell & 2 \\
Code sauber und nachvollziehbar & 1 \\
\midrule
\textbf{Total} & \textbf{15} \\
\bottomrule
\end{tabular}
\caption{Bewertungsraster Transferaufgabe}
\label{tab:bewertung-transfer-15}
\end{table}

\section{Präsentation}

\begin{myexercise}[Kurzvortrag 3 Minuten]
Präsentieren Sie Ihr Projekt in \textbf{3 Minuten} anhand folgender Struktur:
\begin{enumerate}
    \item Problemstellung und Datensatz (30 Sek.)
    \item Gewählte Modelle und Ergebnisse (90 Sek.)
    \item Wichtigste Erkenntnis / Reflexion (60 Sek.)
\end{enumerate}
Verwenden Sie maximal \textbf{2 Folien oder Plots} als visuelle Unterstützung.
\end{myexercise}

\begin{myremark}[Hilfreiche Ressourcen für eigene Datensätze]
\begin{itemize}
    \item \href{https://www.kaggle.com/datasets}{Kaggle Datasets}
    \item \href{https://archive.ics.uci.edu/}{UCI Machine Learning Repository}
    \item \href{https://opendata.swiss/de}{opendata.swiss} (Schweizer Open Data)
\end{itemize}
\end{myremark}
